\itemhead{复习题二 · 第12题（教材p.62）}{XJ-B1-C02-RV-P062-Y0994}
\label{XJ-B1-C02-RV-P062-Y0994}
{\footnotesize 湘教版教材p.62 / PDF第69页 / 复习题二 / 第12题}
\field{原题}{设二次函数 $y=kx^2-kx+\dfrac34$。
\subq{1}若方程 $y=0$ 有实根，则实数 $k$ 的取值范围是\blank。
\subq{2}若不等式 $y>0$ 的解集为 $\varnothing$，则实数 $k$ 的取值范围是\blank。
\subq{3}若不等式 $y>0$ 的解集为 $\mathbb R$，则实数 $k$ 的取值范围是\blank。}
\field{参考答案}{题目限定二次函数，故 $k\ne0$。（1）$(-\infty,0)\cup[3,+\infty)$；（2）$\varnothing$，因为 $y(0)=3/4>0$；（3）$(0,3)$。若取消二次函数限定，第三问才可另含 $k=0$。}
{\small \field{关联知识}{I06 恒成立、有解与范围内最值（教材帮p.71、72、73 / PDF80、81、82）；Q01 二次函数、零点、判别式与根系关系（教材帮p.65 / PDF74）}
\field{方法与易错点}{先明确主元及参数，分降次情况，再用判别式、配方或区间最值。 易错：严格号与非严格号混淆；把有解条件当恒成立；直接用全域判别式处理有限区间。}
\field{答案依据}{答案册相关题组：PDF第15、16页（册内页14、15）；按当前题库原题逐项重算，不按旧版题号机械对套。}}
